\scanpage{51}
\mspage{42}
% archive sheet number: 51

i.e. la \unsure{sempiternelle} condition
\[
X_j^\circ\cap X_i\ne\varnothing\Longrightarrow j\leq i.
\]
Mais on a
\[
X_j^\circ\cap X_i\subset X_j\cap X_i=\bigcup_{k\leq i,j}X_k
\]
et si on n'avait pas $j\leq i$, on ne peut \unsure{avec} la somme du dernier mbr avoir $k=j$, donc cette somme est $\subset\bigcup_{k<j}X_k=\partial X_j\subset\complement X_j^\circ$ \inserted{$\subset\complement X_j^\circ\cap X_i$} donc $X_j^\circ\cap X_i$ est contenu dans son complémentaire, ce qui contredit l'hyp. que c'est vide.

La prop. est prouvée.

\uuline{Corollaire} $3$ Pour que $(X_i)_{i\in I}$ soit une figure \unsure{monoblaste} indexée par l'ens ordonné $I$ (cf Cor. $2$), il f et il s. qu'on ait les conditions suivantes
\[
\forall i,j\in I\ \text{avec }i\leq j,\ \text{on a }X_i\deleted{\unsure{\in}}\subset X_j\tag*{a)}
\]
\[
\forall i,j\in I,\ \text{on a }\deleted{X_i^\circ\ne\varnothing,\ \unsure{\Longleftrightarrow i\ne j}}\quad\inserted{i\ne j\Longrightarrow}\ X_i^\circ\cap X_j^\circ=\varnothing\tag*{\unsure{b)}}
\]
% slanted
\margin{\uline{NB} b équivaut à $2$ conditions $b')\ \forall i\in I,\ X_i^\circ\ne\varnothing$ $b'')\ \forall i,j,\ i\ne j\quad X_i^\circ\cap X_j^\circ=\varnothing$}

\unsure{c)} $\forall i\in I$, si on désigne par $\unsure{\Phi|P}_i\subset\unsure{\mathcal{P}|\mathrm{P}}(X)$ l'ens des parties de $X$ de la forme $X_j$, $j\leq i$, $\Phi_i$ est une figure \unsure{monoblaste}.

\deleted{on \unsure{trouve|prouvera}} \deleted{(\unsure{d'après} \illegible{2})}
\inserted{(Il suffit $=$ qu'elle soit ``quasi-admissible'', et \unsure{moins|même} encore, que \deleted{\unsure{$\forall$|tout}} $x\in X_i$, dans l'ens des $X_j$ ($j\leq i$) qui contiennent $x$, qu'il y ait un élément minimal (ce qui est \unsure{certain} p.ex. $I$ est fini !!))}

Tout revient à montrer que ces conditions \deleted{a) b)} impliquent \deleted{\unsure{pour le moins que}} \inserted{$3')\quad X_i=\bigcup_{j\leq i}X_j^\circ.$} (car on \deleted{\unsure{l'a}} \deleted{$\forall i,j\in I$} \unsure{a} \deleted{$2')$, et} $2')=b'')$.
\[
\deleted{X_i\subset X_j\Longleftrightarrow i\leq j}\qquad\deleted{(\Longleftarrow\text{ est connu par a})}
\]
% slanted
\margin{\uline{NB} les conditions b) \deleted{\unsure{c)}} impliquent déjà (c) via $\bigcup X_i^\circ=$}

\deleted{Ce qui impliquera : $2')$ et $3')$ donc (i).}

\deleted{Prouvons, $\forall i\in I$}
\[
\begin{aligned}
\deleted{X_i^\circ}&\deleted{=\text{cellule \unsure{ouverte|associée} \unsure{de|à} }X_i\text{ dans }\Phi_i.}\\
&\deleted{=\text{(on a donc) }X_i\setminus\bigcup_{\substack{j\leq i\\X_j\subsetneq X_i}}X_j}
\end{aligned}
\]
\[
\deleted{\unsure{\subset}}
\]
